% GENERATED FILE. Edit canonical lessons, then run npm run generate:books.
% canonical-source-hash: fnv1a64:b35a85825a2490c5
% canonical-lessons: FR-C467-tetu, FR-C467-paresse, FR-C467-avare, FR-C467-maladroit, FR-C467-etourdi

\chapter{Qualités et défauts}
\label{ch:fr-qualites-et-defauts}
\hlchaptermodality{467}

\begin{chapteropening}
\textbf{By the end of this chapter:} \emph{I can describe the weak sides of someone's character.}

Complete the last lesson of chapter 467: I can describe the weak sides of someone's character.
\end{chapteropening}

\section[têtu]{têtu — stubborn}

\label{lesson:FR-C467-tetu}



\begin{quote}
Before the new one: say the French for to sulk, then the French for to tease.
\end{quote}

\subsection*{You'll want to know: têtu}
\textbf{têtu} — \textquotedblleft{}stubborn\textquotedblright{}.

\begin{quote}
Mon grand-père est têtu : il ne change jamais d'avis.
\end{quote}

\begin{quote}
\emph{My grandfather is stubborn: he never changes his mind.}
\end{quote}

Built on \emph{la tête}: someone with a hard head. Feminine \emph{têtue}.

\begin{cousinweb}
\emph{têtu} is built on Old French \emph{teste} \textquotedblleft{}head\textquotedblright{} (EtymDB), now \emph{tête}, which is inherited from Latin \emph{testa} \textquotedblleft{}an earthen pot, a shell\textquotedblright{} (Lewis \& Short).
\end{cousinweb}

\begin{grammarlens}[title={What you've built}]
One more word for telling, explaining and describing at B1.
\end{grammarlens}

\subsection*{Your turn}
\begin{itemize}
\raggedright
  \item \emph{Say it:} \emph{têtu}
  \item \emph{Say it:} \emph{têtu}, once more
  \item \emph{Say it:} a sentence of your own with \emph{têtu}
  \item \emph{Recall:} say the French for to sulk, then the French for to tease, then say \emph{têtu} again
\end{itemize}

\begin{tcolorbox}[breakable,colback=teal!4,colframe=teal!35!black,title={Before you move on}]
What does têtu mean? (\textquotedblleft{}Stubborn\textquotedblright{}.) Say it once more.
\end{tcolorbox}
\section[la paresse]{la paresse — laziness}

\label{lesson:FR-C467-paresse}



\begin{quote}
Before the new one: say the French for to tease, then the French for stubborn.
\end{quote}

\subsection*{You'll want to know: la paresse}
\textbf{la paresse} — \textquotedblleft{}laziness\textquotedblright{}.

\begin{quote}
Le dimanche, je me laisse aller à la paresse.
\end{quote}

\begin{quote}
\emph{On Sundays I give in to laziness.}
\end{quote}

The noun of \emph{paresseux} \textquotedblleft{}lazy\textquotedblright{}.

\begin{cousinweb}
\emph{la paresse} is inherited, through Old French \emph{peresce}, from Latin \emph{pigritia} \textquotedblleft{}sluggishness\textquotedblright{} (EtymDB), built on \emph{piger} \textquotedblleft{}reluctant, slow\textquotedblright{} (Lewis \& Short).
\end{cousinweb}

\begin{grammarlens}[title={What you've built}]
One more word for telling, explaining and describing at B1.
\end{grammarlens}

\subsection*{Your turn}
\begin{itemize}
\raggedright
  \item \emph{Say it:} \emph{la paresse}
  \item \emph{Say it:} \emph{la paresse}, once more
  \item \emph{Say it:} a sentence of your own with \emph{la paresse}
  \item \emph{Recall:} say the French for to tease, then the French for stubborn, then say \emph{la paresse} again
\end{itemize}

\begin{tcolorbox}[breakable,colback=teal!4,colframe=teal!35!black,title={Before you move on}]
What does la paresse mean? (\textquotedblleft{}Laziness\textquotedblright{}.) Say it once more.
\end{tcolorbox}
\section[avare]{avare — stingy, mean with money; miser}

\label{lesson:FR-C467-avare}



\begin{quote}
Before the new one: say the French for stubborn, then the French for laziness.
\end{quote}

\subsection*{You'll want to know: avare}
\textbf{avare} — \textquotedblleft{}stingy, mean with money; miser\textquotedblright{}.

\begin{quote}
Il n'est pas avare, mais il compte chaque euro.
\end{quote}

\begin{quote}
\emph{He is not stingy, but he counts every euro.}
\end{quote}

Also a noun: \emph{un avare} — a miser.

\begin{cousinweb}
\emph{avare} is a learned borrowing from Latin \emph{avārus} \textquotedblleft{}greedy\textquotedblright{}, related to \emph{avēre} \textquotedblleft{}to long for, to crave\textquotedblright{} (Lewis \& Short).
\end{cousinweb}

\begin{grammarlens}[title={What you've built}]
One more word for telling, explaining and describing at B1.
\end{grammarlens}

\subsection*{Your turn}
\begin{itemize}
\raggedright
  \item \emph{Say it:} \emph{avare}
  \item \emph{Say it:} \emph{avare}, once more
  \item \emph{Say it:} a sentence of your own with \emph{avare}
  \item \emph{Recall:} say the French for stubborn, then the French for laziness, then say \emph{avare} again
\end{itemize}

\begin{tcolorbox}[breakable,colback=teal!4,colframe=teal!35!black,title={Before you move on}]
What does avare mean? (\textquotedblleft{}Stingy, mean with money; miser\textquotedblright{}.) Say it once more.
\end{tcolorbox}
\section[maladroit]{maladroit — clumsy}

\label{lesson:FR-C467-maladroit}



\begin{quote}
Before the new one: say the French for laziness, then the French for stingy.
\end{quote}

\subsection*{You'll want to know: maladroit}
\textbf{maladroit} — \textquotedblleft{}clumsy\textquotedblright{}.

\begin{quote}
Je suis maladroit : j'ai encore cassé un verre.
\end{quote}

\begin{quote}
\emph{I am clumsy: I broke another glass.}
\end{quote}

The opposite is \emph{adroit} \textquotedblleft{}skilful\textquotedblright{}. \emph{Une maladresse} is a clumsy act.

\begin{cousinweb}
\emph{maladroit} is built from \emph{mal} \textquotedblleft{}badly\textquotedblright{} and \emph{adroit} \textquotedblleft{}skilful\textquotedblright{} (EtymDB). \emph{Adroit} comes from \emph{à droit} \textquotedblleft{}rightly\textquotedblright{}, and \emph{droit} is inherited from Latin \emph{dīrēctus} \textquotedblleft{}straight\textquotedblright{}.
\end{cousinweb}

\begin{grammarlens}[title={What you've built}]
One more word for telling, explaining and describing at B1.
\end{grammarlens}

\subsection*{Your turn}
\begin{itemize}
\raggedright
  \item \emph{Say it:} \emph{maladroit}
  \item \emph{Say it:} \emph{maladroit}, once more
  \item \emph{Say it:} a sentence of your own with \emph{maladroit}
  \item \emph{Recall:} say the French for laziness, then the French for stingy, then say \emph{maladroit} again
\end{itemize}

\begin{tcolorbox}[breakable,colback=teal!4,colframe=teal!35!black,title={Before you move on}]
What does maladroit mean? (\textquotedblleft{}Clumsy\textquotedblright{}.) Say it once more.
\end{tcolorbox}
\section[étourdi]{étourdi — scatterbrained, absent-minded}

\label{lesson:FR-C467-etourdi}



\begin{quote}
Before the new one: say the French for stingy, then the French for clumsy.
\end{quote}

\subsection*{You'll want to know: étourdi}
\textbf{étourdi} — \textquotedblleft{}scatterbrained, absent-minded\textquotedblright{}.

\begin{quote}
Ma sœur est un peu étourdie ; elle oublie souvent ses clés.
\end{quote}

\begin{quote}
\emph{My sister is a bit scatterbrained; she often forgets her keys.}
\end{quote}

Also \emph{dazed}, after a blow: \emph{il est encore étourdi}.

\begin{cousinweb}
\emph{étourdi} is the participle of \emph{étourdir} \textquotedblleft{}to daze\textquotedblright{}. It is usually traced to a Vulgar Latin \emph{exturdīre}, built on \emph{turdus} \textquotedblleft{}a thrush\textquotedblright{}, a bird once thought to be easily dazed; the origin is not fully settled.
\end{cousinweb}

\begin{grammarlens}[title={What you've built}]
One more word for telling, explaining and describing at B1.
\end{grammarlens}

\subsection*{Your turn}
\begin{itemize}
\raggedright
  \item \emph{Say it:} \emph{étourdi}
  \item \emph{Say it:} \emph{étourdi}, once more
  \item \emph{Say it:} a sentence of your own with \emph{étourdi}
  \item \emph{Recall:} say the French for stingy, then the French for clumsy, then say \emph{étourdi} again
\end{itemize}

\begin{tcolorbox}[breakable,colback=teal!4,colframe=teal!35!black,title={Before you move on}]
What does étourdi mean? (\textquotedblleft{}Scatterbrained, absent-minded\textquotedblright{}.) Say it once more.
\end{tcolorbox}
